% Part: first-order-logic
% Chapter: sequent-calculus
% Section: rules-and-proofs

\documentclass[../../../include/open-logic-section]{subfiles}

\begin{document}

\olfileid{mvl}{seq}{rul}

\olsection{قاعدې او \usetoken{P}{derivation}}

په لاندې بحث کښې دې~$\Gamma, \Delta, \Pi, \Lambda$ د
!!{sentence}s متناهي لړۍ تمثيل کړي۔

\begin{defn}[سېکوېنټ]
يو \emph{$n$ اړخيز سېکوېنټ} د لاندې بڼې عبارت دے:
\[
\Gamma_1 \nSequent \dots \nSequent \Gamma_n
\]
چې هره~$\Gamma_i$ د ژبې~$\Lang L$ د !!{sentence}s يوه متناهي
(کېدے شي تشه) لړۍ ده۔ د سرچينې يادونه: په اصلي متن کښې دلته
د هرې لړۍ دپاره د لومړي شاخص نښه ليکل شوې؛ د ټولو اړخونو دپاره
ئې عمومي شاخص ته واړوله۔
\end{defn}

\begin{defn}[ابتدايي سېکوېنټ]
يو \emph{$n$ اړخيز ابتدايي سېکوېنټ} د هرې ژبنۍ !!{sentence}~$!A$
دپاره د~$!A \nSequent \dots \nSequent !A$ په بڼه~$n$ اړخيز
سېکوېنټ دے۔

که ژبه يو~$0$ ځايه نښلوونکے~$\star$، يعنې بياني ثابت، ولري، نو
هغه سېکوېنټ~$\dots \nSequent \star \nSequent \dots$ هم ابتدايي
ګڼو چې~$\star$ پکې د~$\tf{\star} \in V$ له اړوند رښتياوالي
قيمت سره په برابر ځاے کښې راځي او نور ځايونه ئې تش وي۔
\end{defn}

د~$n$ ارزښته منطق~$\Log{L}$ د هر نښلوونکي دپاره، په~$\Log{L}$ کښې
د هغه د هر ممکن رښتياوالي قيمت يوه منطقي قاعده شته۔ د~$\Log{L}$ د~$n$ اړخيز
سېکوېنټ حساب !!^{derivation}s د سېکوېنټونو ونې دي: تر ټولو بره
سېکوېنټونه ابتدايي وي، او که يو سېکوېنټ د يوه يا څو نورو لاندې
ولاړ وي، نو بايد د~$\Log{L}$ د نښلوونکو د استنتاج د يوې قاعدې له
مخې په سمه توګه ترې راوځي۔

\begin{defn}[قضيې]
!!^a{sentence}~$!A$ د~$n$ ارزښته منطق~$\Log{L}$ \emph{قضيه}
ده که د هغه~$n$ اړخيز سېکوېنټ !!a{derivation} وي چې د~$\Log{L}$
د هر ټاکل شوي رښتياوالي قيمت له اړوند ځاے سره~$!A$ لري۔ که~$!A$
قضيه وي،~$\Proves[\Log{L}] !A$ ليکو، او که نۀ وي،~$\Proves/[\Log{L}]
!A$ ليکو۔
\end{defn}

\begin{defn}[!!^{derivability}]
!!^a{sentence}~$!A$ د~$n$ ارزښته منطق~$\Log{L}$ کښې د
!!{sentence}s له سټ~$\Gamma$ څخه \emph{!!{derivable} ده}، يعنې~$\Gamma
\Proves[\Log{L}] !A$، هله او يوازې هله چې يو متناهي
فرعي سټ~$\Gamma_0 \subseteq \Gamma$ او د~$\Gamma_0$ د
!!{sentence}s يوه لړۍ~$\Gamma_0'$ وي چې لاندې سېکوېنټ
!!a{derivation} ولري:
\[ \Lambda_1 \nSequent \dots \nSequent \Lambda_n \]
دلته~$\Lambda_i$ هغه مهال~$!A$ دے چې~$i$ ځاے له ټاکل شوي
رښتياوالي قيمت سره برابر وي، او په بل حالت کښې~$\Gamma_0'$ دے۔
که~$!A$ له~$\Gamma$ څخه !!{derivable} نۀ وي،~$\Gamma \Proves/ !A$
ليکو۔
\end{defn}

د بېلګې په توګه، د لوکاشېوېچ~$3$ ارزښته منطق يو~$3$ اړخيز
سېکوېنټ حساب لري۔ په~$3$ اړخيز سېکوېنټ~$\Gamma \nSequent \Pi
\nSequent \Delta$ کښې~$\Gamma$ له~$\False$،~$\Delta$ له~$\True$
او~$\Pi$ له~$\Undef$ سره برابر دي۔ بديهي سېکوېنټونه~$!A
\nSequent !A \nSequent !A$ دي۔ ځکه چې يوازې~$\True$ ټاکل شوے دے،
$\Gamma \Proves[\LogLuk[3]] !A$ هله او يوازې هله چې سېکوېنټ~$\Gamma
\nSequent \Gamma \nSequent !A$ يو~!!a{derivation} ولري۔ (که~$\Undef$
هم ټاکل شوے وے، نو د~$\Gamma \nSequent !A \nSequent !A$
!!a{derivation} به پکار و۔)

\end{document}
